\newif\ifglobal

%%%%%%%%%%%%%%%%%%%%%%%
%%% Compile to view %%%
%%%%%%%%%%%%%%%%%%%%%%%

%%%%%%%%%%%%%%%%%%%%%%%%%%%%%%%%%%%%%%%%%%%%%%%%%%%
%%% Readme:                                     %%%
%%% https://www.overleaf.com/read/bwdjvfjksvjy/ %%%
%%%%%%%%%%%%%%%%%%%%%%%%%%%%%%%%%%%%%%%%%%%%%%%%%%%

% >>> M?: marks a module setting number or important notice

% >>> M4: Set {./relative-path} to external files for this module <<<
\def \modulefiles {./30-AES}
% To add an external file, use:
% (Replace '---' with actual filename)
% '\input{\modulefiles/tables/---.tex}' for tables
% '\includegraphics{\modulefiles/figures/---}' for figures
% '\subfile{\modulefiles/---}' for register files

% >>> M1: This module is in global project? <<<
% 'YES' - uncomment; 'NO' - comment out
\globaltrue

\ifglobal
    \documentclass[main\_\_CN.tex]{subfiles}

\else
    % Font "FandolSong-Regular" used by ctex does not have GSUB table.
    % It causes warnings that do not affect compilation and can be ignored.
    \PassOptionsToPackage{quiet}{fontspec}

    \documentclass[openany, 10pt]{book}

    % >>> M2: In ./00-shared/preamble-trm-module, also find
    %         settings for visibility of todo notes, watermark <<<
    \usepackage{preamble-trm-module}


    % Enable Chinese typesetting
    \usepackage{ctex}

    % >>> M3: Temporary labels in Chinese
    %         you can add your own labels there <<<
    \myexternaldocument{./00-shared/temp-labels-cn}

\fi %\ifglobal


\setcounter{secnumdepth}{4}

% >>> M5: If this module needs any updates in
% './00-shared/preamble-trm-module.sty'
% (include additional package, etc.), contact your module PM
% or create the file './00-shared/changelog.tex' make the
% necessary updates yourself and document them in this file <<<


\begin{document}


% Sets language into which pre-defined bits of text are translated
\selectlanguage{chinese}

% Do not delete this block
\ifglobal
% Add the following front matter if
% this module is outside of global project
\else
    \InputIfFileExists{readme}{}{}
    {\let\clearpage\relax \listoftodos}
    \tableofcontents
    %\listoftables
    %\listoffigures
\fi


%%%%%%%%%%%%%%%%%%%%%%%%%
%%% TRM Module Begins %%%
%%%%%%%%%%%%%%%%%%%%%%%%%

\hypertarget{aes}{}
\chapter{AES 加速器 (AES)}\label{mod:aes}

\section{概述}

\chipname{} 内置 AES（高级加密标准）硬件加速器可使用 AES 算法，完成数据的加解密运算，具有 \hyperref[sec:aes-typical-modes]{Typical AES} 和 \hyperref[sec:aes-dma-aes]{DMA-AES} 两种工作模式。整体而言，相比基于纯软件的 AES 运算，AES 硬件加速器能够极大地提高运算速度。

\section{主要特性}

\chipname{} 支持以下特性：
\begin{itemize}
    \item Typical AES 工作模式
        \begin{itemize}
            \item AES-128/AES-256 加解密运算
        \end{itemize}
    \end{itemize}

    \begin{itemize}
    \item DMA-AES 工作模式
        \begin{itemize}
            \item AES-128/AES-256 加解密运算
            \item 块（加密）模式
                \begin{itemize}
                    \item ECB (Electronic Codebook)
                    \item CBC (Cipher Block Chaining)
                    \item OFB (Output Feedback)
                    \item CTR (Counter)
                    \item CFB8 (8-bit Cipher Feedback)
                    \item CFB128 (128-bit Cipher Feedback)
                \end{itemize}
            \item 中断发生
        \end{itemize}
    \end{itemize}


\section{工作模式简介}

\chipname{} 内置的 AES 加速器支持 Typical AES 和 DMA-AES 两种工作模式。

\begin{itemize}
    \item Typical AES 工作模式：
    \begin{itemize}
        \item 支持使用 128 位或 256 位密钥进行加密与解密运算，即 \href{https://nvlpubs.nist.gov/nistpubs/FIPS/NIST.FIPS.197.pdf}{NIST FIPS 197} 标准中的 AES-128 和 AES-256  加解密运算。
    \end{itemize}
    这种情况下，明文/密文的读/写操作统一通过 CPU 访问完成。
    \item DMA-AES 工作模式：
    \begin{itemize}
        \item 支持使用 128 位或 256 位密钥进行加密与解密运算，即 \href{https://nvlpubs.nist.gov/nistpubs/FIPS/NIST.FIPS.197.pdf}{NIST FIPS 197} 标准中的 AES-128 和 AES-256 加解密运算；
        \item 还支持 \href{https://nvlpubs.nist.gov/nistpubs/Legacy/SP/nistspecialpublication800-38a.pdf}{NIST SP 800-38A} 标准中的 ECB/CBC/OFB/CTR/CFB8/CFB128 等块加密模式运算。
    \end{itemize}
    在这种情况下，明文/密文的传输通过硬件上的 DMA 完成，计算完成时会有中断发生。
\end{itemize}

用户可通过配置 \hyperref[regdesc:AESDMAENABLEREG]{AES\_DMA\_ENABLE\_REG} 选择 AES 加速器的工作模式，具体参考表 \ref{tab:aes-dma-enable}。

\begin{longtable}[c]{ |c|l| }
\caption{工作模式}
\label{tab:aes-dma-enable}
\\\hline
\rowcolor{lightgray}
\hyperref[regdesc:AESDMAENABLEREG]{AES\_DMA\_ENABLE\_REG}       & 工作模式\\
    \hline
    0           & Typical AES    \\
    \hline
    1           & DMA-AES        \\
    \hline
\end{longtable}

用户可通过配置 \hyperref[regdesc:AESMODEREG]{AES\_MODE\_REG} 寄存器选择密钥长度和加解密方向，具体可参考表 \ref{tab:aes-operation}。

\begin{longtable}[c]{ | c | l | }
\caption{密钥长度和加解密方向}
\label{tab:aes-operation}
\\\hline
\rowcolor{lightgray}
\hyperref[regdesc:AESMODEREG]{AES\_MODE\_REG}[2:0]     & 密钥长度和加解密方向\\
    \hline
    0                   & AES-128 加密    \\
    \hline
    1                   & 保留    \\
    \hline
    2                   & AES-256 加密    \\
    \hline
    3                   & 保留    \\
    \hline
    4                   & AES-128 解密    \\
    \hline
    5                   & 保留    \\
    \hline
    6                   & AES-256 解密    \\
    \hline
    7                   & 保留    \\
    \hline
\end{longtable}

有关 Typical AES 和 DMA-AES 两种工作模式的具体介绍，请见下方 \ref{sec:aes-typical-modes} 章节和 \ref{sec:aes-dma-aes} 章节。

\begin{importantlisting}
\vspace{-0.5em}

    \chipname{} 的 \nameref{mod:digsig} 模块也会调用 AES 加速器。此时，用户无法正常访问 AES 加速器。
\end{importantlisting}

\section{Typical AES 工作模式}\label{sec:aes-typical-modes}

在 Typical AES 工作模式下，
AES 加速器的状态值可查看寄存器 \hyperref[regdesc:AESSTATEREG]{AES\_STATE\_REG}，具体见表 \ref{tab:aes-state-value-typical} 所示：

\begin{longtable}[c]{ |c | l | l | }
\caption{状态返回值}
\label{tab:aes-state-value-typical}
\\\hline
\rowcolor{lightgray}
返回值 & 描述 &  状态说明 \\ \hline

    0                   & IDLE      &   加速器空闲或计算完成\\ \hline
    1                   & WORK      &   加速器忙于计算     \\ \hline
\end{longtable}

\subsection{密钥、明文、密文}

寄存器 \hyperref[regdesc:AESKEYnREG]{AES\_KEY\_\regindex{n}\_REG} 用于存放密钥，由 8 个 32 位寄存器组成。

\begin{itemize}
    \item 如果为 AES-128 加解密运算，则 128 位密钥在寄存器 \hyperref[regdesc:AESKEYnREG]{AES\_KEY\_0\_REG} \verb+~+ \hyperref[regdesc:AESKEYnREG]{AES\_KEY\_3\_REG} 中。
    \item 如果为 AES-256 加解密运算，则 256 位密钥在寄存器 \hyperref[regdesc:AESKEYnREG]{AES\_KEY\_0\_REG} \verb+~+ \hyperref[regdesc:AESKEYnREG]{AES\_KEY\_7\_REG} 中。
\end{itemize}

寄存器 \hyperref[regdesc:AESTEXTINnREG]{AES\_TEXT\_IN\_\regindex{m}\_REG} 和 \hyperref[regdesc:AESTEXTOUTnREG]{AES\_TEXT\_OUT\_\regindex{m}\_REG} 用于存放明文和密文，各由 4 个 32 位寄存器组成。

\begin{itemize}
    \item 如果为 AES-128/265 加密运算，则运算开始之前用明文初始化寄存器 \hyperref[regdesc:AESTEXTINnREG]{AES\_TEXT\_IN\_\regindex{m}\_REG}。运算完成之后，AES 加速器将把密文更新入寄存器 \hyperref[regdesc:AESTEXTOUTnREG]{AES\_TEXT\_OUT\_\regindex{m}\_REG}。
    \item 如果为 AES-128/256 解密运算，则运算开始之前用密文初始化寄存器 \hyperref[regdesc:AESTEXTINnREG]{AES\_TEXT\_IN\_\regindex{m}\_REG}。运算完成之后，AES 加速器将把明文更新入寄存器 \hyperref[regdesc:AESTEXTOUTnREG]{AES\_TEXT\_OUT\_\regindex{m}\_REG}。
\end{itemize}

\subsection{字节序}

\textbf{文本字节序}

在 Typical AES 工作模式下，AES 加速器可以使用密钥对 128 位的 block 进行加解密。在操作寄存器 \hyperref[regdesc:AESTEXTINnREG]{AES\_TEXT\_IN\_\regindex{m}\_REG} 和 \hyperref[regdesc:AESTEXTINnREG]{AES\_TEXT\_OUT\_\regindex{m}\_REG} 中的数据时，用户应遵循表 \ref{tab:aes-endianness-type-typical-plaintext-ciphertext} 中定义的文本字节序。


\begin{table}[H]
    \centering
    \caption{Typical AES 文本字节序}
    \label{tab:aes-endianness-type-typical-plaintext-ciphertext}

    \begin{threeparttable}

    \begin{tabular}{ |c|c|c|c|c|c|}
    \hline
        \rowcolor{lightgray}
        \multicolumn{6}{|c|}{明文/密文} \\\hline
        \multicolumn{2}{|c|}{\multirow{2}{*}{State\tnote{1}}} &\multicolumn{4}{c|}{c\tnote{2}}\\ \cline{3-6}
        \multicolumn{2}{|c|}{} &0 &1 &2 &3\\ \cline{1-6}
        \multirow{4} {*}r & 0 &\small{\hyperref[regdesc:AESTEXTINnREG]{AES\_TEXT\_\regindex{x}\_0\_REG}[7:0]} &   \small{\hyperref[regdesc:AESTEXTINnREG]{AES\_TEXT\_\regindex{x}\_1\_REG}[7:0]}  &   \small{\hyperref[regdesc:AESTEXTINnREG]{AES\_TEXT\_\regindex{x}\_2\_REG}[7:0]}  &   \small{\hyperref[regdesc:AESTEXTINnREG]{AES\_TEXT\_\regindex{x}\_3\_REG}[7:0]}\\
        \cline{2-6}
        &1       &   \small{\hyperref[regdesc:AESTEXTINnREG]{AES\_TEXT\_\regindex{x}\_0\_REG}[15:8]} &   \small{\hyperref[regdesc:AESTEXTINnREG]{AES\_TEXT\_\regindex{x}\_1\_REG}[15:8]}     &   \small{\hyperref[regdesc:AESTEXTINnREG]{AES\_TEXT\_\regindex{x}\_2\_REG}[15:8]} &   \small{\hyperref[regdesc:AESTEXTINnREG]{AES\_TEXT\_\regindex{x}\_3\_REG}[15:8]}\\
        \cline{2-6}
        &2   &   \small{\hyperref[regdesc:AESTEXTINnREG]{AES\_TEXT\_\regindex{x}\_0\_REG}[23:16]}    &   \small{\hyperref[regdesc:AESTEXTINnREG]{AES\_TEXT\_\regindex{x}\_1\_REG}[23:16]}    &   \small{\hyperref[regdesc:AESTEXTINnREG]{AES\_TEXT\_\regindex{x}\_2\_REG}[23:16]}    &   \small{\hyperref[regdesc:AESTEXTINnREG]{AES\_TEXT\_\regindex{x}\_3\_REG}[23:16]}\\
        \cline{2-6}
        &3   &   \small{\hyperref[regdesc:AESTEXTINnREG]{AES\_TEXT\_\regindex{x}\_0\_REG}[31:24]}    &   \small{\hyperref[regdesc:AESTEXTINnREG]{AES\_TEXT\_\regindex{x}\_1\_REG}[31:24]}    & \small{\hyperref[regdesc:AESTEXTINnREG]{AES\_TEXT\_\regindex{x}\_2\_REG}[31:24]}  &   \small{\hyperref[regdesc:AESTEXTINnREG]{AES\_TEXT\_\regindex{x}\_3\_REG}[31:24]}\\
        \hline
        \end{tabular}

        \begin{tablenotes}
            \item[1] 有关 ``State（以及 c 和 r）" 的详细定义，请参考 \href{https://nvlpubs.nist.gov/nistpubs/FIPS/NIST.FIPS.197.pdf}{NIST FIPS 197} 中 ``3.4 The State" 章节。
            \item[2] 其中，\regindex{x} = IN 或 OUT。
        \end{tablenotes}
     \end{threeparttable}
    \end{table}

\textbf{密钥字节序}

在 Typical AES 工作模式下，在向寄存器 \hyperref[regdesc:AESKEYnREG]{AES\_KEY\_\regindex{n}\_REG} 中填入数据时，用户应遵循表 \ref{tab:aes-128-key-endian} 和表 \ref{tab:aes-256-key-endian} 中定义的文本字节序。


\begin{table}[H]
    \centering
    \caption{AES-128 密钥字节序}
    \label{tab:aes-128-key-endian}
    \begin{threeparttable}

    \begin{tabular}{ |l|l|l|l|l| }
    \hline
        \rowcolor{lightgray}
        Bit\tnote{1} &   w[0]    &   w[1]    &   w[2]    &   w[3]\tnote{2}   \\\hline
        [31:24]     &   \hyperref[regdesc:AESKEYnREG]{AES\_KEY\_0\_REG}[7:0]    &   \hyperref[regdesc:AESKEYnREG]{AES\_KEY\_1\_REG}[7:0]    &   \hyperref[regdesc:AESKEYnREG]{AES\_KEY\_2\_REG}[7:0]    &   \hyperref[regdesc:AESKEYnREG]{AES\_KEY\_3\_REG}[7:0]\\ \hline
        [23:16]     &   \hyperref[regdesc:AESKEYnREG]{AES\_KEY\_0\_REG}[15:8]   &   \hyperref[regdesc:AESKEYnREG]{AES\_KEY\_1\_REG}[15:8]   &   \hyperref[regdesc:AESKEYnREG]{AES\_KEY\_2\_REG}[15:8]   &   \hyperref[regdesc:AESKEYnREG]{AES\_KEY\_3\_REG}[15:8]\\ \hline
        [15:8]  &   \hyperref[regdesc:AESKEYnREG]{AES\_KEY\_0\_REG}[23:16]  &   \hyperref[regdesc:AESKEYnREG]{AES\_KEY\_1\_REG}[23:16]  &   \hyperref[regdesc:AESKEYnREG]{AES\_KEY\_2\_REG}[23:16]  &   \hyperref[regdesc:AESKEYnREG]{AES\_KEY\_3\_REG}[23:16]\\ \hline
        [7:0]   &   \hyperref[regdesc:AESKEYnREG]{AES\_KEY\_0\_REG}[31:24]  &   \hyperref[regdesc:AESKEYnREG]{AES\_KEY\_1\_REG}[31:24]  & \hyperref[regdesc:AESKEYnREG]{AES\_KEY\_2\_REG}[31:24]    &   \hyperref[regdesc:AESKEYnREG]{AES\_KEY\_3\_REG}[31:24]\\ \hline
        \end{tabular}

        \begin{tablenotes}
            \item[1] Bit 列代表 w[0] \verb+~+ w[3] 每个 word 中的各个字节。
            \item[2] w[0] \verb+~+ w[3] 符合标准 \href{https://nvlpubs.nist.gov/nistpubs/FIPS/NIST.FIPS.197.pdf}{NIST FIPS 197} 中 ``5.2 Key Expansion" 章节中对 ``the first Nk words of the expanded key" 的描述。
        \end{tablenotes}
    \end{threeparttable}
\end{table}

\newgeometry{left=1cm,right=1cm,top=2cm,bottom=2cm}
\begin{landscape}

\begin{threeparttable}

\begin{longtable}[c]
{ |l|l|l|l|l|l|l|l|l|}
\caption{AES-256 密钥字节序}\label{tab:aes-256-key-endian}\\
\hline
\rowcolor{lightgray}
Bit$^1$ &   w[0]    &   w[1]    &   w[2]    &   w[3]    &   w[4]    &   w[5]&   w[6]    &   w[7]$^2$    \\\hline
\endfirsthead
\hline
\rowcolor{lightgray}
Bit &   w[0]    &   w[1]    &   w[2]    &   w[3]    &   w[4]    &   w[5]&   w[6]    &   w[7]    \\\hline
\endhead
\endfoot
[31:24]     &   \scriptsize \hyperref[regdesc:AESKEYnREG]{AES\_KEY\_0\_REG}[7:0]    &   \scriptsize \hyperref[regdesc:AESKEYnREG]{AES\_KEY\_1\_REG}[7:0]    &   \scriptsize \hyperref[regdesc:AESKEYnREG]{AES\_KEY\_2\_REG}[7:0]    &   \scriptsize \hyperref[regdesc:AESKEYnREG]{AES\_KEY\_3\_REG}[7:0]    &   \scriptsize \hyperref[regdesc:AESKEYnREG]{AES\_KEY\_4\_REG}[7:0]    &   \scriptsize \hyperref[regdesc:AESKEYnREG]{AES\_KEY\_5\_REG}[7:0]
& \scriptsize  \hyperref[regdesc:AESKEYnREG]{AES\_KEY\_6\_REG}[7:0]    &   \scriptsize \hyperref[regdesc:AESKEYnREG]{AES\_KEY\_7\_REG}[7:0]\\ \hline
[23:16]     &   \scriptsize \hyperref[regdesc:AESKEYnREG]{AES\_KEY\_0\_REG}[15:8]   &   \scriptsize \hyperref[regdesc:AESKEYnREG]{AES\_KEY\_1\_REG}[15:8]   &   \scriptsize \hyperref[regdesc:AESKEYnREG]{AES\_KEY\_2\_REG}[15:8]   &   \scriptsize \hyperref[regdesc:AESKEYnREG]{AES\_KEY\_3\_REG}[15:8]   &   \scriptsize \hyperref[regdesc:AESKEYnREG]{AES\_KEY\_4\_REG}[15:8]   &   \scriptsize \hyperref[regdesc:AESKEYnREG]{AES\_KEY\_5\_REG}[15:8]   &   \scriptsize \hyperref[regdesc:AESKEYnREG]{AES\_KEY\_6\_REG}[15:8]   &   \scriptsize \hyperref[regdesc:AESKEYnREG]{AES\_KEY\_7\_REG}[15:8]\\ \hline
[15:8]  &  \scriptsize \hyperref[regdesc:AESKEYnREG]{AES\_KEY\_0\_REG}[23:16]  &   \scriptsize \hyperref[regdesc:AESKEYnREG]{AES\_KEY\_1\_REG}[23:16]  &   \scriptsize \hyperref[regdesc:AESKEYnREG]{AES\_KEY\_2\_REG}[23:16]  &   \scriptsize \hyperref[regdesc:AESKEYnREG]{AES\_KEY\_3\_REG}[23:16]  &   \scriptsize \hyperref[regdesc:AESKEYnREG]{AES\_KEY\_4\_REG}[23:16]  &   \scriptsize \hyperref[regdesc:AESKEYnREG]{AES\_KEY\_5\_REG}[23:16]& \scriptsize \hyperref[regdesc:AESKEYnREG]{AES\_KEY\_6\_REG}[23:16]  &   \scriptsize \hyperref[regdesc:AESKEYnREG]{AES\_KEY\_7\_REG}[23:16]\\ \hline
[7:0]   &  \scriptsize \hyperref[regdesc:AESKEYnREG]{AES\_KEY\_0\_REG}[31:24]  &    \scriptsize \hyperref[regdesc:AESKEYnREG]{AES\_KEY\_1\_REG}[31:24]  & \scriptsize \hyperref[regdesc:AESKEYnREG]{AES\_KEY\_2\_REG}[31:24]    &   \scriptsize \hyperref[regdesc:AESKEYnREG]{AES\_KEY\_3\_REG}[31:24]  &   \scriptsize \hyperref[regdesc:AESKEYnREG]{AES\_KEY\_4\_REG}[31:24]  &   \scriptsize \hyperref[regdesc:AESKEYnREG]{AES\_KEY\_5\_REG}[31:24]& \scriptsize \hyperref[regdesc:AESKEYnREG]{AES\_KEY\_6\_REG}[31:24]  &   \scriptsize \hyperref[regdesc:AESKEYnREG]{AES\_KEY\_7\_REG}[31:24]\\ \hline
\end{longtable}

\begin{tablenotes}
\item[1] Bit 列代表 w[0] \verb+~+ w[7] 每个 word 中的各个字节。
\item[2] w[0] \verb+~+ w[7] 符合标准 \href{https://nvlpubs.nist.gov/nistpubs/FIPS/NIST.FIPS.197.pdf}{NIST FIPS 197} 中 ``5.2 Key Expansion" 章节中对 ``the first Nk words of the expanded key" 的描述。
\end{tablenotes}
\end{threeparttable}

\end{landscape}
\restoregeometry

\subsection{Typical AES 工作模式的流程 }

\textbf{单次运算}

\begin{enumerate}
    \item 对寄存器 \hyperref[regdesc:AESDMAENABLEREG]{AES\_DMA\_ENABLE\_REG} 写入 0。
    \item 初始化寄存器 \hyperref[regdesc:AESMODEREG]{AES\_MODE\_REG}、\hyperref[regdesc:AESKEYnREG]{AES\_KEY\_\regindex{n}\_REG}、\hyperref[regdesc:AESTEXTINnREG]{AES\_TEXT\_IN\_\regindex{m}\_REG}。
    \item 启动运算。对寄存器 \hyperref[regdesc:AESTRIGGERREG]{AES\_TRIGGER\_REG} 写入 1。
    \item 等待运算完成。轮询寄存器 \hyperref[regdesc:AESSTATEREG]{AES\_STATE\_REG}，直到读到 0。
    \item 从寄存器 \hyperref[regdesc:AESTEXTOUTnREG]{AES\_TEXT\_OUT\_\regindex{m}\_REG} 读取结果。
\end{enumerate}

\textbf{连续运算}

在连续运算过程中，每次运算完成之后，只有寄存器 \hyperref[regdesc:AESTEXTINnREG]{AES\_TEXT\_IN\_\regindex{m}\_REG} 和 \hyperref[regdesc:AESTEXTOUTnREG]{AES\_TEXT\_OUT\_\regindex{m}\_REG} (\regindex{m}: 0-3) 会被 AES 加速器更新，
而 \hyperref[regdesc:AESDMAENABLEREG]{AES\_DMA\_ENABLE\_REG}、\hyperref[regdesc:AESMODEREG]{AES\_MODE\_REG}、\hyperref[regdesc:AESKEYnREG]{AES\_KEY\_\regindex{n}\_REG} 等寄存器中的内容不会变化。
所以进行连续运算时可以简化初始化操作。

\begin{enumerate}
    \item 第一次运算之前对寄存器 \hyperref[regdesc:AESDMAENABLEREG]{AES\_DMA\_ENABLE\_REG} 写入 0。
    \item 第一次运算之前初始化寄存器 \hyperref[regdesc:AESMODEREG]{AES\_MODE\_REG} 和 \hyperref[regdesc:AESKEYnREG]{AES\_KEY\_\regindex{n}\_REG}。
    \item 更新寄存器 \hyperref[regdesc:AESTEXTINnREG]{AES\_TEXT\_IN\_\regindex{m}\_REG}。 \label{others:aes-update-register-typical-continuous}
    \item 启动运算。对寄存器 \hyperref[regdesc:AESTRIGGERREG]{AES\_TRIGGER\_REG} 写入 1。
    \item 等待运算完成。轮询寄存器 \hyperref[regdesc:AESSTATEREG]{AES\_STATE\_REG}，直到读到 0。
    \item 从寄存器 \hyperref[regdesc:AESTEXTOUTnREG]{AES\_TEXT\_OUT\_\regindex{m}\_REG} 读取结果 \label{read result in typical AES continuous calculation}。 返回步骤 \ref{others:aes-update-register-typical-continuous}，进行下一轮运算。
\end{enumerate}


\clearpage

\section{DMA-AES 工作模式}\label{sec:aes-dma-aes}

在 DMA-AES 工作模式下，AES 加速器可支持 ECB/CBC/OFB/CTR/CFB8/CFB128 等 6 种块模式运算。
用户可通过配置 \hyperref[regdesc:AESBLOCKMODEREG]{AES\_BLOCK\_MODE\_REG} 寄存器选择具体运算类型，具体可参考表 \ref{tab:aes-block-mode}。

\begin{longtable}[c]{ | c | l | }
\caption{块模式选择}
\label{tab:aes-block-mode}
\\\hline
\rowcolor{lightgray}
\hyperref[regdesc:AESBLOCKMODEREG]{AES\_BLOCK\_MODE\_REG}[2:0]     & 块模式\\
    \hline
    0                          & ECB (Electronic Code Book)    \\
    \hline
    1                          & CBC (Cipher Block Chaining)     \\
    \hline
    2                          & OFB (Output FeedBack)    \\
    \hline
    3                          & CTR (Counter)    \\
    \hline
    4                          & CFB8 (8-bit Cipher FeedBack)    \\
    \hline
    5                          & CFB128 (128-bit Cipher FeedBack) \\
    \hline
    6                          & 保留     \\
    \hline
    7                          & 保留     \\
    \hline
\end{longtable}

AES 加速器的状态值可查看寄存器 \hyperref[regdesc:AESSTATEREG]{AES\_STATE\_REG}，具体见表 \ref{tab:aes-state-value-dma} 所示：

\begin{longtable}[c]{ | c | l | l | }
\caption{状态返回值}
\label{tab:aes-state-value-dma}
\\\hline
\rowcolor{lightgray}
返回值 & 描述 &  状态说明 \\ \hline

    0                   & IDLE      &   加速器空闲\\ \hline
    1                   & WORK      &   加速器忙于计算     \\ \hline
    2                   & DONE      &   加速器计算完成     \\ \hline
\end{longtable}

AES 加速器在 DMA-AES 工作模式下允许中断发生，软件清零。
中断功能默认关闭，用户可通过将 \hyperref[regdesc:AESINTENAREG]{AES\_INT\_ENA\_REG} 寄存器配置为 1 开启中断。如开启中断功能，AES 加速器在完成计算时，中断发生。

\subsection{密钥、明文、密文}

\textbf{块运算模式}

在块运算模式下，AES 加速器的源数据来自 DMA，结果数据也将被写入 DMA。
\begin{itemize}
    \item 如果为加密运算，则 DMA 从 memory 中读取明文数据流并将其传给 AES。AES 计算出密文后将密文写入 DMA。DMA 再将密文写入 memory。
    \item 如果为解密运算，则 DMA 从 memory 中读取密文数据流并将其传给 AES。AES 计算出明文后将明文写入 DMA。DMA 再将明文写入 memory。
\end{itemize}

AES 加速器在进行块运算时，结果数据与源数据的大小保持一致。此时，DMA 的数据搬运过程和 AES 的计算过程有所交叠，因此总工作时间有所减少。

\label{api:testpadding} 值得注意的是，AES 加速器在 DMA-AES 工作模式下要求源数据的大小必须是 128 位的整数倍，否则需要将原始明文封装为 128 位的整数倍，即在原比特串 (bit string) 尾部尽可能少的补 “0”，具体过程见表 \ref{tab:aes-define-func-text-padding} 所示。

\clearpage
\begin{longtable}[c]{|p{15cm}|}
\caption{TEXT-PADDING} 
\label{tab:aes-define-func-text-padding}\\
\hline
\rowcolor{lightgray}
\multirow{1}{*}{%
{\bfseries Function :} {\bfseries TEXT-PADDING}(\ ) } \\\hline
{\bfseries Input\phantom{aaa} :} $X$, bit string.\\
{\bfseries Output\phantom{a~} :} $Y$ = {\bfseries TEXT-PADDING}($X$), whose length is the nearest integral multiples of 128 bits. \\\hline
{\bfseries Steps\phantom{aa~}  }\\
\hspace{1.3cm} Let us assume that X is a data-stream that can be split into {\it{n}} parts as following: \\
\phantom{aaaaaaaa}$X = X_{1} || X_{2} || \cdots || X_{n-1} || X_{n}$ \\
\phantom{aaaaaaaa}Here, the lengths of $X_{1}, X_{2}, \cdots, X_{n-1}$ all equal to 128 bits, and the length of $X_{n}$ is {\it{t}} (0<={\it{t}}<=127).\\
\phantom{aaaaaaaa}If ${\it{t}} = 0$, then\\
\phantom{aaaaaaaa}\phantom{空空空空空} {\bfseries TEXT-PADDING}($X$) = $X$;\\
\phantom{aaaaaaaa}If $0<{\it{t}} <= 127$, define a 128-bit block, $X_{n}^{*}$, and let $X_{n}^{*} = X_{n} || 0^{128-{\it{t}}}$, then\\
\phantom{aaaaaaaa}\phantom{空空空空空} {\bfseries TEXT-PADDING}($X$) = $X_{1} || X_{2} || \cdots || X_{n-1} || X_{n}^{*}$  = $X || 0^{128-{\it{t}}}$\\
\hline
\end{longtable}


\subsection{字节序}

在 DMA-AES 工作模式下，源数据和结果数据的传输完全由 DMA 完成，因此不支持字节序的控制调节，但要求它们在 memory 中以一定的方式来存放，且要求数据量必须是 block 的整数倍。

举例说明，假设 DMA 需要搬运 2 个 block 大小的源数据：

\begin{itemize}
    \item 十六进制： 0102030405060708090A0B0C0D0E0F101112131415161718191A1B1C1D1E1F20
\end{itemize}

假设起始地址为 0x{}0280，则源数据在 memory 中的存放位置如表 \ref{tab:aes-dma-aes-plaintext-ciphertext-endian} 所示。结果数据也遵从相同的存放规则，在此不多做介绍。

\begin{longtable}[c]
{ |c|c|c|c|c|c|c|c|}
\caption{DMA AES 存储字节序} 
\label{tab:aes-dma-aes-plaintext-ciphertext-endian}\\
\hline
\rowcolor{lightgray}
地址 &  字节  & 地址 &  字节  & 地址 &  字节  & 地址 &  字节 \\\hline
\endfirsthead
\hline
\rowcolor{lightgray}
地址 &  字节  & 地址 &  字节  & 地址 &  字节  & 地址 &  字节 \\\hline
\endhead
\endfoot
0x{}0280 & 0x{}01 & 0x{}0281 & 0x{}02 & 0x{}0282 & 0x{}03 & 0x{}0283 & 0x{}04 \\\hline
0x{}0284 & 0x{}05 & 0x{}0285 & 0x{}06 & 0x{}0286 & 0x{}07 & 0x{}0287 & 0x{}08 \\\hline
0x{}0288 & 0x{}09 & 0x{}0289 & 0x{}0A & 0x{}028A & 0x{}0B & 0x{}028B & 0x{}0C \\\hline
0x{}028C & 0x{}0D & 0x{}028D & 0x{}0E & 0x{}028E & 0x{}0F & 0x{}028F & 0x{}10 \\\hline
0x{}0290 & 0x{}11 & 0x{}0291 & 0x{}12 & 0x{}0292 & 0x{}13 & 0x{}0293 & 0x{}14 \\\hline
0x{}0294 & 0x{}15 & 0x{}0295 & 0x{}16 & 0x{}0296 & 0x{}17 & 0x{}0297 & 0x{}18 \\\hline
0x{}0298 & 0x{}19 & 0x{}0299 & 0x{}1A & 0x{}029A & 0x{}1B & 0x{}029B & 0x{}1C \\\hline
0x{}029C & 0x{}1D & 0x{}029D & 0x{}1E & 0x{}029E & 0x{}1F & 0x{}029F & 0x{}20 \\\hline
\end{longtable}

另外，值得注意的是，DMA 既可以访问片内存储空间，又可以访问片外 PSRAM。当访问片外 PSRAM 时，基地址必须满足 DMA 对地址的相关要求，但当访问片内存储器空间时则没有限制。详情请见章节 \ref{mod:dma} \textit{\nameref{mod:dma}}。


\subsection{标准增量函数}\label{sec:aes-Standard-Incrementing-Function}

AES 加速器在进行 CTR 块运算时，还可提供两种标准增量函数供用户选择：INC$_{32}$ 和 INC$_{128}$ 。用户可通过将寄存器 \hyperref[regdesc:AESINCSELREG]{AES\_INC\_SEL\_REG} 置为 0 或 1 选择 INC$_{32}$ 或 INC$_{128}$ 标准增量函数。更多有关标准增量函数的内容，请见 \href{https://nvlpubs.nist.gov/nistpubs/Legacy/SP/nistspecialpublication800-38a.pdf}{NIST SP 800-38A} 标准中的 “B.1 The Standard Incrementing Function” 章节。


\subsection{块个数}\label{sec:aes-block-number}

寄存器 \hyperref[regdesc:AESBLOCKNUMREG]{AES\_BLOCK\_NUM\_REG} 存放明文或密文的块个数 (Block Number)，其值等于 length({\bfseries TEXT-PADDING}($P$))/128，也等于 length({\bfseries TEXT-PADDING}($C$))/128。这里的 $P$ 指明文 (plaintext)， $C$ 指密文 (ciphertext)。该寄存器仅在 DMA-AES 工作模式下有意义。

\subsection{初始向量}\label{sec:aes-init-vector}

\label{memdesc:AESIVMEM} 存储器 \hyperref[memdesc:AESIVMEM]{AES\_IV\_MEM} 的空间大小为 16 字节，仅在块运算模式下有效。
对于 CBC/OFB/CFB8/CFB128 等操作，\hyperref[memdesc:AESIVMEM]{AES\_IV\_MEM} 用于存放初始向量 (Initialization Vector, IV) 的值。对于 CTR 操作， \hyperref[memdesc:AESIVMEM]{AES\_IV\_MEM}  存放初始计数器 (Initial Counter Block, ICB) 的值。

IV 和 ICB 都是 128-bit 长的比特串，从左向右被分割成 16 个字节 (Byte0, Byte1, Byte2, \begin{math}\cdots\end{math}, Byte15)， 构成一个字节序列，在 \hyperref[memdesc:AESIVMEM]{AES\_IV\_MEM} 中存放时需要遵循表 \ref{tab:aes-dma-aes-plaintext-ciphertext-endian} 中的字节序规则，即 Byte0 存放在 \hyperref[memdesc:AESIVMEM]{AES\_IV\_MEM} 中的最低地址中，Byte15 存放在 \hyperref[memdesc:AESIVMEM]{AES\_IV\_MEM} 中的最高地址中。

更多有关 IV 和 ICB 的信息，请参考 \href{https://nvlpubs.nist.gov/nistpubs/Legacy/SP/nistspecialpublication800-38a.pdf}{NIST SP 800-38A} 标准。

\subsection{DMA-AES 工作模式的流程 }

\begin{enumerate}
    \item 选择一条 DMA 通道与 AES 加速器连接，配置 DMA 链表，而后启动 DMA。详情请见章节 \ref{mod:dma} \textit{\nameref{mod:dma}}。
    \item 配置 AES：
    \begin{itemize}
        \item 对寄存器 \hyperref[regdesc:AESDMAENABLEREG]{AES\_DMA\_ENABLE\_REG} 写入 1。
        \item 选择是否开启中断。根据需要设置寄存器 \hyperref[regdesc:AESINTENAREG]{AES\_INT\_ENA\_REG} 的值。
        \item 初始化 \hyperref[regdesc:AESMODEREG]{AES\_MODE\_REG} 和 \hyperref[regdesc:AESKEYnREG]{AES\_KEY\_\regindex{n}\_REG} 寄存器。
        \item 配置 \hyperref[regdesc:AESBLOCKMODEREG]{AES\_BLOCK\_MODE\_REG} 寄存器，选择具体块加密模式。详见表 \ref{tab:aes-block-mode}。
        \item 初始化寄存器 \hyperref[regdesc:AESBLOCKNUMREG]{AES\_BLOCK\_NUM\_REG}，请参照章节 \ref{sec:aes-block-number}。
        \item 初始化寄存器 \hyperref[regdesc:AESINCSELREG]{AES\_INC\_SEL\_REG}（仅在 CTR 块模式下使用）。
        \item 初始化存储器 \hyperref[memdesc:AESIVMEM]{AES\_IV\_MEM}（在 ECB 块模式下不使用）。
    \end{itemize}
    \item 启动运算。对寄存器 \hyperref[regdesc:AESTRIGGERREG]{AES\_TRIGGER\_REG} 写入 1。
    \item \label{wait ECB/CBC/OFB/CTR/CFB8/CFB128 done}
            等待运算完成。轮询寄存器 \hyperref[regdesc:AESSTATEREG]{AES\_STATE\_REG}，直到读到 2。
            如果开启了中断功能，也可以等待 AES\_INT 中断产生。
    \item 确认 DMA 完成从 AES 到内存的数据传输。此时，结果数据已经被 DMA 写入 memory，可以直接从中读取。详情请参考章节 \ref{mod:dma} \textit{\nameref{mod:dma}}。
    \item 如果开启了中断，当处理中断程序完成后，请及时对寄存器 \hyperref[regdesc:AESINTCLRREG]{AES\_INT\_CLR\_REG} 写 1 以清除中断。
    \item 对寄存器 \hyperref[regdesc:AESDMAEXITREG]{AES\_DMA\_EXIT\_REG} 写入 1 释放 AES 加速器。之后如果再读取寄存器 \hyperref[regdesc:AESSTATEREG]{AES\_STATE\_REG} 将读到 0。
    该步操作可以提前完成，但必须在步骤 \ref{wait ECB/CBC/OFB/CTR/CFB8/CFB128 done} 之后。
\end{enumerate}

\clearpage
\section{存储器列表}

本小节的所有地址均为相对于 AES 加速器基地址的地址偏移量（相对地址），具体基地址请见章节 \ref{mod:sysmem} \textit{\nameref{mod:sysmem}} 中的表 \ref{tab:sysmem-base-address}。

\subfile{\modulefiles/30-AES-mem-summ__CN}

\clearpage
\section{寄存器列表}\hypertarget{aes-reg-summ}{}

本小节的所有地址均为相对于 AES 加速器基地址的地址偏移量（相对地址），具体基地址请见章节 \ref{mod:sysmem} \textit{\nameref{mod:sysmem}} 中的表 \ref{tab:sysmem-base-address}。

请查看章节 \hyperref[glossary-access-types]{\textit{寄存器的访问类型}}，了解“访问”列缩写的含义。

\subfile{\modulefiles/30-AES-reg-summ__CN}


\clearpage
\section{寄存器}

本小节的所有地址均为相对于 AES 加速器基地址的地址偏移量（相对地址），具体基地址请见章节 \ref{mod:sysmem} \textit{\nameref{mod:sysmem}} 中的表 \ref{tab:sysmem-base-address}。

\subfile{\modulefiles/30-AES-reg__CN}

\end{document}
